\subsection{The first middle attachment}\label{sec:attachment}

The filtration has been constructed without a nonvanishing
assumption.  We now ask whether its first height-one layer attaches
nontrivially.  The answer depends on the actual canonical
\(K(2)\)-localization map, not on the rational dimensions of
the layers.

First note why rationalizing the filtration does not answer this
question.  Its rationalized cofibers are single \(Q\)-modules
in the strictly decreasing degrees
\(0,-1,-3,-4,-5,-6,-8,-9\).
Inductively, if the previous rationalized boundaries vanish,
\(L_0F_{i-1}\) is the sum of the earlier shifts of \(Q\).
The degree of the new cofiber is smaller than every earlier
degree, and hence cannot be a homotopy degree of
\(\Sigma L_0F_{i-1}\).  A \(Q\)-module map from the new shift
to that suspension is zero.  This proves by induction that
every rationalized boundary vanishes.  A spectral boundary can
nevertheless be nonzero.

\subsubsection{A conditional lifting criterion}

Let
\[
 u_X:L_0X\longrightarrow L_0L_{K(2)}X
\]
be induced by the canonical localization unit.  Retain the
specific \(h:U_{\mathrm{mid}}\to Y\), \(\beta_3\), and
\(d_3=\partial_wh|_{U_3}\) constructed in
\Cref{sec:application}, where \(U_3=\Sigma^{-3}L_1S\).

\begin{proposition}[Canonical-map criterion]\label{prop:attachment-criterion}
If \(u_X\) does not kill
\(\beta_3\in\pi_{-3}L_0X\), then
\[
 d_3:U_3=\Sigma^{-3}L_1S\longrightarrow\Sigma W
\]
is nonzero as a map of underlying spectra.
\end{proposition}

\begin{proof}
Suppose \(d_3\) were null after forgetting the module structure.
Applying spectral mapping spaces from \(U_3\) to the
cofiber sequence
\[
 W\longrightarrow X\xrightarrow qY\xrightarrow{\partial_w}\Sigma W
\]
shows that the chosen nullhomotopy lifts \(h|_{U_3}\) to a map
\(\widetilde h_3:U_3\to X\).
Rationally, \(L_0W\) has homotopy only in degrees \(0\) and
\(-1\).  Its groups in degrees \(-3\) and \(-4\) are zero,
so the cofiber sequence makes
\(\pi_{-3}L_0X\to\pi_{-3}L_0Y\) an isomorphism.
The construction of \(h\) sends the degree-\(-3\) generator
of \(L_0U_3\) to the projection of \(\beta_3\).
Thus \(L_0\widetilde h_3\) sends it to \(\beta_3\) itself.
This uses only the image of that underlying homotopy element.

But \(L_{K(2)}U_3=0\), since \(U_3\) is \(E(1)\)-local.
By localization adjunction, every map from \(U_3\) to the
\(K(2)\)-local target \(L_{K(2)}X\) is null.  In particular,
the composite of \(\widetilde h_3\) with the unit
\(X\to L_{K(2)}X\) is null.  Rationalizing and using naturality
would make \(u_X\) kill \(\beta_3\), contrary to the hypothesis.
\end{proof}

This criterion is independent of the existence argument:
its antecedent was not used to construct \(h\), the stages, or
their terminal equivalence.

\subsubsection{The companion canonical-map theorem}

For every prime \(p\geq5\), the canonical map for the uncompleted
sphere \(\mathbb S\),
\begin{equation}\label{eq:companion-map}
 L_0L_{K(3)}\mathbb S
 \longrightarrow L_0L_{K(2)}L_{K(3)}\mathbb S,
\end{equation}
is nonzero on \(\pi_{-3}\).
This is \cite[Theorem~1.1]{CompanionHeightThreeObstruction}.

\begin{corollary}[Nonzero first middle attachment]\label{cor:nonzero-attachment}
For every prime \(p\geq5\), retain the common height-three choice
in \Cref{sec:application}, which realizes \Cref{thm:main} and the
height-three instance of \Cref{thm:attachment-identification}.
Then
\[
 d_3:\Sigma^{-3}L_1S\longrightarrow
       \Sigma\bigl(L_2S\vee\Sigma^{-1}L_2S\bigr)
\]
is nonzero as a map of underlying spectra.
\end{corollary}

\begin{proof}
Set \(T_0=L_{K(3)}\mathbb S\).  The completion map
\(\mathbb S\to S\) is a mod-\(p\) equivalence.
After \(p\)-localization its cofiber has invertible
multiplication by \(p\), hence is rational and \(K(3)\)-acyclic.
Completion therefore induces an equivalence
\(e_c:T_0\xrightarrow{\simeq}T\).
Naturality of the \(K(2)\)-localization unit gives
\begin{equation}\label{eq:completion-square}
\begin{tikzcd}[column sep=large]
 L_0T_0 \arrow[r] \arrow[d,"L_0e_c"',"\simeq"]
   & L_0L_{K(2)}T_0
       \arrow[d,"L_0L_{K(2)}e_c","\simeq"']\\
 L_0T \arrow[r] & L_0L_{K(2)}T.
\end{tikzcd}
\end{equation}
The top map is nonzero on \(\pi_{-3}\) by
\cite[Theorem~1.1]{CompanionHeightThreeObstruction}.  Hence so is the bottom map.

Let \(\ell:T\to X=L_2T\) be the localization map.
Its cofiber is \(E(2)\)-acyclic.  Since
\(\langle E(2)\rangle=\langle K(0)\vee K(1)\vee K(2)\rangle\),
that cofiber is both rationally and \(K(2)\)-acyclic.
The vertical maps in the next naturality square are therefore
equivalences:
\begin{equation}\label{eq:localization-square}
\begin{tikzcd}[column sep=large]
 L_0T \arrow[r] \arrow[d,"L_0\ell"',"\simeq"]
   & L_0L_{K(2)}T
       \arrow[d,"L_0L_{K(2)}\ell","\simeq"']\\
 L_0X \arrow[r,"u_X"] & L_0L_{K(2)}X.
\end{tikzcd}
\end{equation}
Thus \(u_X\) is nonzero on \(\pi_{-3}\), and it is the same
canonical map under these identifications.  An arbitrary nonzero
map between the same two objects would not suffice.

The square \eqref{eq:localization-square} is a square of
\(S\)-modules by \eqref{eq:module-localization}.
Rationalization is extension to \(Q=H\Qp\), so its homotopy
maps are \(\Qp\)-linear.  Equivalently, the action of
\(\pi_0S=\Zp\) on rational homotopy extends uniquely across
inversion of \(p\).  By \Cref{h3:thm:rational-compatibility},
\[
 \pi_{-3}L_0X=\Qp\alpha_3.
\]
A nonzero \(\Qp\)-linear map from this one-dimensional space
cannot kill \(\alpha_3\) or its nonzero scalar multiple
\(\beta_3=c^{-1}\alpha_3\).  This conclusion requires no
numerical identification of a primitive chosen in the companion
proof with the chosen \(\alpha_3\).

The hypothesis of \Cref{prop:attachment-criterion} is now
satisfied, so \(d_3\ne0\).
\end{proof}

Thus the companion theorem enters only in verifying the antecedent of
the conditional criterion.  The stages and local bases in
\Cref{thm:main} were constructed independently.
